\documentclass[12pt]{article}
\usepackage{amsmath}
\usepackage{geometry}
\geometry{margin=1in}

\title{CSE400 -- Fundamentals of Probability in Computing}
\author{Instructor: Dhaval Patel, PhD}
\date{March 10, 2026}

\begin{document}

\maketitle

\section*{Lecture L15--L17: Joint Random Variables + Joint Distributions + Examples}

\section*{1. Outline}

\begin{itemize}
    \item Why Study Joint Random Variables?
    \item Discrete Case
    \begin{itemize}
        \item Joint PMF Table Representation
        \item Example: Rolling Two Dice
    \end{itemize}
    \item Continuous Case
    \begin{itemize}
        \item Geometric Interpretation
        \item Worked Example
    \end{itemize}
    \item Joint Cumulative Distribution Function
    \begin{itemize}
        \item Definition of Joint CDF
        \item Continuous Case
        \item Discrete Case
        \item Properties of the Joint CDF
    \end{itemize}
    \item Comprehensive Example
\end{itemize}

\section*{2. Why Study Joint Random Variables?}

\subsection*{Example \#1: Smart Transportation}

Two random variables:
\begin{itemize}
    \item $X =$ Vehicle Speed
    \item $Y =$ Traffic Density
\end{itemize}

\textbf{Question:} Can speed be high when traffic density is also high?\\
\textbf{Answer:} Yes / No

\[
X \text{ and } Y \text{ are correlated random variables}
\]

\textbf{Applications:}
\begin{itemize}
    \item Adaptive traffic signals
    \item Congestion prediction
    \item Autonomous vehicle routing
\end{itemize}

\subsection*{Example \#2: Wireless Network Performance}

Two random variables:
\begin{itemize}
    \item $X =$ Signal-to-Noise Ratio (SNR)
    \item $Y =$ Data Throughput
\end{itemize}

\textbf{Why Joint RVs?}\\
Throughput depends on SNR (both vary due to fading, interference, congestion)

\textbf{Question:} If SNR is high, will throughput always be high?\\
\textbf{Answer:} Yes / No

\textbf{Reason:} Network congestion (conditional relationships)

\textbf{Applications:}
\begin{itemize}
    \item 5G scheduling
    \item Adaptive modulation
    \item Network planning
\end{itemize}

\subsection*{General Observation}

In science and real life, we are often interested in more than one random variable simultaneously.

\subsection*{Example \#3: Weather Prediction}

\begin{itemize}
    \item $X =$ Rainfall
    \item $Y =$ Temperature
\end{itemize}

Application: Weather forecasting

\subsection*{Example \#4: Drone Delivery}

\begin{itemize}
    \item $X =$ Wind Speed
    \item $Y =$ Delivery Time
\end{itemize}

Application: Logistics and UAV systems

\subsection*{Example \#5: Rolling Two Dice}

\begin{itemize}
    \item $(X, Y) =$ outcomes of dice
\end{itemize}

Conditions:
\begin{itemize}
    \item $X + Y = 7$
    \item $X > Y$
\end{itemize}

Requires joint PMF

\section*{3. Discrete Case -- Joint PMF}

(Title slides introducing concept)

\section*{4. Joint PMF Table Representation}

Joint PMF is represented using a table:
\begin{itemize}
    \item Rows $\rightarrow$ values of one random variable
    \item Columns $\rightarrow$ values of another random variable
    \item Each cell $\rightarrow$ joint probability
\end{itemize}

\section*{5. Example: Rolling Two Dice}

\begin{itemize}
    \item Two dice rolled
    \item Define:
    \begin{itemize}
        \item $X =$ outcome of first die
        \item $Y =$ outcome of second die
    \end{itemize}
\end{itemize}

Joint PMF table constructed for all pairs $(x, y)$

(All possible combinations considered)

\section*{6. Continuous Case -- Joint PDF}

(Series of slides introducing concept)

\subsection*{Joint PDF (Continuous Case)}

\begin{itemize}
    \item Used when random variables are continuous
    \item Defined over a region in the plane
    \item Represents density instead of probability
\end{itemize}

\subsection*{Geometric Interpretation}

\begin{itemize}
    \item Probability corresponds to area under surface
    \item Region in $(x, y)$-plane defines probability
\end{itemize}

\section*{7. Example (Continuous Case)}

(Problem introduced)

\subsection*{Example Solution}

(Steps shown across slides)

\begin{itemize}
    \item Define region
    \item Integrate joint PDF over region
    \item Compute required probability
\end{itemize}

\section*{8. Joint Cumulative Distribution Function (Joint CDF)}

\subsection*{Definition}

Joint CDF:
\[
F_{X,Y}(x,y) = P(X \le x, Y \le y)
\]

\section*{9. Joint CDF -- Continuous Case}

\begin{itemize}
    \item Computed using double integration over region:
    \begin{itemize}
        \item From $-\infty$ to $x$
        \item From $-\infty$ to $y$
    \end{itemize}
\end{itemize}

\section*{10. Joint CDF -- Discrete Case}

\begin{itemize}
    \item Computed using summation:
    \begin{itemize}
        \item Over all values satisfying:
        \[
        X \le x, \quad Y \le y
        \]
    \end{itemize}
\end{itemize}

\section*{11. Properties of the Joint CDF}

\begin{itemize}
    \item Non-decreasing function
    \item Limits:
    \begin{itemize}
        \item $F_{X,Y}(-\infty, -\infty) = 0$
        \item $F_{X,Y}(\infty, \infty) = 1$
    \end{itemize}
    \item Right-continuous
    \item Defines probabilities over regions
\end{itemize}

\section*{12. End of Lecture}

\textbf{Thank You}

\end{document}